\chapter{System Architecture}
\label{ch:architecture}

\section{Overall Architecture}
\label{sec:overall_arch}

The system is organized as a five-stage pipeline. Raw business events enter at the left and
a structured Business Intelligence Briefing exits at the right, ready for consumption by any
LLM agent. Figure~\ref{fig:system_overview} shows the full pipeline with annotated token
counts at each stage transition.

\begin{figure}[H]
  \centering
  \includegraphics[width=\textwidth]{system_overview}
  \caption{Five-stage neuro-symbolic distillation pipeline. Token counts are illustrative
           for a 1,000-event batch.}
  \label{fig:system_overview}
\end{figure}

The \textbf{Data Generator} produces synthetic supply chain events at a configurable rate.
Each event carries a type, a timestamp, a product identifier, a supplier or warehouse
identifier, and a set of numerical attributes (quantities, prices, defect rates, ETAs).
In a production deployment, this stage would be replaced by a connector to an ERP system
or a message queue.

The \textbf{Graph Builder} processes each event batch and updates the knowledge graph in
memory. Node attributes are updated to reflect the latest state of each business entity.
New nodes and edges are added when new relationships appear in the event stream. The graph
is snapshotted before and after each batch so that the delta engine can compare the two
states.

The \textbf{Delta Engine} takes the before and after snapshots and computes a structural
diff. The diff is represented as five categories of change: new nodes, removed nodes,
changed node attributes, new edges, and removed edges. Each category of change is
individually scoreable and has distinct business semantics---a new edge between a supplier
and a product means a new sourcing relationship, while a change in a node's
\texttt{anomaly\_score} attribute means an escalating problem at that entity.

The \textbf{Saliency Scorer} assigns a scalar importance score to each delta item using
the formula described in Chapter~\ref{ch:methodology}. Delta items below the threshold
$\theta = 0.4$ are classified as noise and discarded. Items above the threshold are
classified as CRITICAL (score $> 0.75$), IMPORTANT ($0.5 < \text{score} \leq 0.75$), or
INFORMATIONAL ($0.4 < \text{score} \leq 0.5$). Only CRITICAL and IMPORTANT items are
included in the briefing forwarded to agents.

The \textbf{Agent Layer} contains three specialized agents---analyst, risk, and forecast---each
receiving the distilled briefing rather than the raw event stream. The orchestrator routes
the briefing to all three agents in parallel and assembles their responses into a composite
report for the dashboard user.

\section{Data Flow Pipeline}
\label{sec:data_flow}

Tracing a single event batch through all stages makes the data flow concrete. Consider a
batch of 1,000 supply chain events covering a 24-hour window of activity.

At the raw input stage, the 1,000 events serialized as JSON occupy roughly 61,000 tokens
when formatted for a language model prompt. Forwarding this directly to an LLM agent would
cost approximately \$1.85 per call at GPT-4o pricing, or \$0.03 per call at Gemini Flash
pricing. For a team running 100 such queries per day, the daily cost is \$185 (GPT-4o) or
\$3 (Gemini Flash).

After the Graph Builder processes the batch, the knowledge graph has absorbed all 1,000
events into a 30-node, 368-edge structure with updated attributes. The graph itself, if
serialized naively, is larger than the event list---but the graph is never serialized for
the prompt. Only the delta computed from it is.

The Delta Engine computes the diff between the pre-batch and post-batch graph snapshots.
A typical 1,000-event batch produces several dozen delta items: a handful of new edges
from new supplier contracts, several changed anomaly scores from quality issues, and a
collection of changed quantities from sales and restocks. The raw delta, before saliency
filtering, occupies roughly 1,500--3,000 tokens if serialized as text.

After saliency scoring and noise filtering, the surviving delta items are assembled into a
structured briefing of 600--900 tokens. This briefing explicitly names each entity
involved, states its current relationship context (``ShenzhenDirect supplies GPU Card RTX
which is stored at WH-LON where 450 units are pending''), and assigns a severity label.
The 84$\times$ token compression at this scale is the ratio of the 61,749-token raw input
to the approximately 734-token average briefing.

\section{Hardware Platform}
\label{sec:hardware}

The system runs on a standard development laptop with no dedicated GPU or specialized
hardware. The software dependencies are: Python 3.12, Streamlit 1.32+, NetworkX 3.2+,
google-generativeai 0.7+ (the Gemini Flash API client), and standard scientific Python
libraries (NumPy, Pandas, Matplotlib). The Streamlit dashboard runs on localhost:8501 and
is accessible from any browser on the same machine.

The only external service dependency is the Gemini Flash API, which requires an API key
and a network connection. The pipeline processing stages (Graph Builder, Delta Engine,
Saliency Scorer, Distiller) run entirely locally and require no internet access. In an
offline or air-gapped deployment, a locally hosted open-source model could replace the
Gemini Flash calls without any change to the pipeline logic.

At the scales tested in this project (up to 3,000 events), the pipeline overhead---the
total time spent in all stages before the LLM call---is under 15 milliseconds. The
bottleneck is the LLM API round-trip, which takes 10--25 seconds depending on the
provider and load. The pipeline overhead contributes less than 0.1\% of end-to-end latency.

\section{Comparison with Naive Baseline}
\label{sec:comparison_baseline}

Table~\ref{tab:method_comparison} compares the three methods evaluated in this project.
The naive baseline forwards the full raw event JSON to the LLM without any preprocessing.
The aggregation baseline computes per-category summary statistics (counts, averages, maxima)
and formats them as a short text summary, discarding individual event details and
relationship context. The full pipeline performs the complete neuro-symbolic distillation
described above.

\begin{table}[H]
\centering
\caption{Comparison of three context preparation methods at 1,000 events}
\label{tab:method_comparison}
\begin{tabular}{lrrrr}
\toprule
\textbf{Method} & \textbf{Tokens} & \textbf{Recall} & \textbf{Latency (ms)} & \textbf{Relative Cost} \\
\midrule
Naive (full context) & 61,749 & 93.8\% & $<$1 & 1.0$\times$ \\
Aggregation only     & 427    & 93.8\% & $<$1 & 0.007$\times$ \\
Full pipeline        & 734    & 93.8\% & 5.6  & 0.012$\times$ \\
\bottomrule
\end{tabular}
\end{table}

The aggregation baseline achieves comparable token counts to the full pipeline and
identical recall on most analytical categories. Its limitation appears on the cascade
detection task: it identifies 1 of 2 affected warehouses in the GT4 cascade pattern,
while the full pipeline identifies both, because the pipeline preserves explicit
relationship context (``ShenzhenDirect $\to$ GPU Card RTX $\to$ WH-LON, WH-NYC'')
that the aggregation baseline's statistical summary discards. This distinction is
discussed in detail in Chapter~\ref{ch:results}.
